\makeatletter
\@ifundefined{keepfig}{\newenvironment{keepfig}{\par\edef\keepfig@pi{\the\parindent}\noindent
  \begin{minipage}{\linewidth}\parindent=\keepfig@pi\relax}{\end{minipage}\par}}{}
\makeatother
\section{Построимые числа ещё раз}

В главе 1 (\rf{1.37}) мы доказали, что вещественное число $\alpha$ построимо циркулем и линейкой тогда и только тогда,
когда оно лежит в подполе $\R$ вида $\Q[\sqrt{a_1},\dots,\sqrt{a_r}]$, где каждое $a_i$~--- положительный элемент поля
$\Q[\sqrt{a_1},\dots,\sqrt{a_{i-1}}]$. Отсюда (\rf{1.38})
\begin{equation}\label{eq:c3-constr7}
\alpha\ \text{построимо}\ \Longrightarrow\ [\Q[\alpha]:\Q]=2^s\ \text{для некоторого}\ s .
\end{equation}
Обратное неверно (\rf{3.27}): степень самого числа~$\alpha$ ничего не говорит о том, как устроены промежуточные поля.
Теория Галуа позволяет доказать частичное обращение \eqref{eq:c3-constr7}, а затем и точный критерий построимости:
важна не степень $\Q[\alpha]$, а степень поля разложения минимального многочлена~$\alpha$.

\begin{theorem}[3.24]
Если $\alpha$ лежит в подполе $E\subset\R$, которое является расширением Галуа поля $\Q$ степени $2^r$, то $\alpha$
построимо.
\end{theorem}

\begin{proof}
Пусть $G=\Gal(E/\Q)$; так как $E/\Q$~--- расширение Галуа, $|G|=[E:\Q]=2^r$ (\rf{3.17}). По лемме \rf{3.24a} ниже (при
$H=\{1\}$) в $G$ есть цепочка подгрупп
\[
\{1\}=G_0\subset G_1\subset G_2\subset\dots\subset G_r=G,\qquad (G_i:G_{i-1})=2 .
\]
По основной теореме (\rf{3.17}) ей отвечает цепочка полей $E_i=E^{G_i}$:
\[
E=E_0\supset E_1\supset E_2\supset\dots\supset E_r=\Q,\qquad [E_{i-1}:E_i]=(G_i:G_{i-1})=2 .
\]
По лемме \rf{3.25} ниже $E_{i-1}=E_i[\sqrt{a_i}]$ для некоторого $a_i\in E_i$; при этом $\sqrt{a_i}\notin E_i$, так что
$a_i\ne0$, и $a_i=(\sqrt{a_i})^2>0$, потому что $\sqrt{a_i}\in E_{i-1}\subset\R$. Итак,
\[
E=\Q\bigl[\sqrt{a_r},\sqrt{a_{r-1}},\dots,\sqrt{a_1}\bigr],\qquad a_i>0,\quad a_i\in E_i=\Q\bigl[\sqrt{a_r},\dots,\sqrt{a_{i+1}}\bigr],
\]
~--- поле в точности того вида, что в \rf{1.37}(b) (с обратной нумерацией). Значит, все элементы $E$, и в частности
$\alpha$, построимы.
\end{proof}

\begin{center}
\begin{keepfig}
\centering
\begin{tikzpicture}[every node/.style={font=\small}]
\node (G0) at (0,0) {$\{1\}=G_0$}; \node (G1) at (0,1.1) {$G_1$}; \node (Gd) at (0,2.2) {$\vdots$};
\node (Gr) at (0,3.3) {$G_r=G$};
\node at (0,0.55) {\rotatebox{90}{$\subset$}}; \node at (0,1.65) {\rotatebox{90}{$\subset$}};
\node at (0,2.75) {\rotatebox{90}{$\subset$}};
\node[cold] at (-1.3,0.55) {$2$}; \node[cold] at (-1.3,1.65) {$2$}; \node[cold] at (-1.3,2.75) {$2$};
\node (E0) at (6,0) {$E=E_0$}; \node (E1) at (6,1.1) {$E_1$}; \node at (6,2.2) {$\vdots$}; \node (Er) at (6,3.3) {$E_r=\Q$};
\node at (6,0.55) {\rotatebox{-90}{$\subset$}}; \node at (6,1.65) {\rotatebox{-90}{$\subset$}};
\node at (6,2.75) {\rotatebox{-90}{$\subset$}};
\node[cnew,right] at (6.5,0.55) {$E_0=E_1[\sqrt{a_1}]$, $a_1>0$};
\node[cnew,right] at (6.5,1.65) {$E_1=E_2[\sqrt{a_2}]$, $a_2>0$};
\node[cnew,right] at (6.5,2.75) {$E_{r-1}=\Q[\sqrt{a_r}]$, $a_r>0$};
\draw[<->,gray] (1.0,0) -- (5.2,0); \draw[<->,gray] (0.5,1.1) -- (5.5,1.1); \draw[<->,gray] (0.9,3.3) -- (5.1,3.3);
\node[gray] at (3,3.55) {$E_i=E^{G_i}$};
\end{tikzpicture}
\figcap{Доказательство \rf{3.24}: цепочке подгрупп $2$-группы с индексами $2$ отвечает башня квадратичных расширений;
каждое из них~--- присоединение квадратного корня (\rf{3.25}), а вещественность $E$ даёт $a_i>0$.}
\end{keepfig}
\end{center}

Отметим, что от группы $G$ здесь нужна не <<разрешимость>> как таковая, а именно цепочка подгрупп, каждая из которых
имеет индекс $2$ в следующей. Подгруппа $G_{i-1}$ нормальна в $G_i$ (подгруппа индекса $2$ всегда нормальна: левый и
правый смежные классы, отличные от неё самой, совпадают с её дополнением), но не обязана быть нормальной в $G$; поэтому
поля $E_i$ не обязаны быть расширениями Галуа поля $\Q$. Это и не нужно: лемма \rf{3.25} применима к любому
квадратичному расширению.

\begin{lemma}[3.24a]
Пусть $G$~--- группа порядка $p^n$, $p$ простое, и $H\subset G$~--- подгруппа. Тогда существует цепочка подгрупп
\[
H=H_0\subset H_1\subset\dots\subset H_k=G,\qquad (H_i:H_{i-1})=p .
\]
В частности (при $H=\{1\}$) существует цепочка $\{1\}=G_0\subset G_1\subset\dots\subset G_n=G$ с $(G_i:G_{i-1})=p$.
\end{lemma}

\begin{proof}
Индукция по $|G|$, а при фиксированной $G$~--- по индексу $(G:H)$. При $H=G$ доказывать нечего. Пусть $H\ne G$ (в
частности, $G\ne\{1\}$). Центр $Z(G)$ нетривиален: в формуле классов $|G|=|Z(G)|+\sum(G:C_G(x))$ (сумма по
нецентральным классам сопряжённости) все индексы $(G:C_G(x))>1$ делят $p^n$ и потому делятся на $p$, так что $p$
делит $|Z(G)|$. Возьмём в $Z(G)$ элемент $z$ порядка $p$ (если $c\in Z(G)$ имеет порядок $p^m$, $m\ge1$, то годится
$z=c^{p^{m-1}}$) и положим $N=\langle z\rangle$; это нормальная подгруппа порядка $p$, все её элементы перестановочны
со всеми элементами $G$.

\emph{Случай $N\not\subset H$.} Тогда $N\cap H=\{1\}$, $HN=NH$~--- подгруппа порядка $p|H|$, и $(HN:H)=p$. Так как
$(G:HN)<(G:H)$, по предположению индукции есть цепочка от $HN$ до $G$ с индексами $p$; добавим к ней в начало $H$.

\emph{Случай $N\subset H$.} Группа $G/N$ имеет порядок $p^{n-1}$, и по предположению индукции для её подгруппы $H/N$
есть цепочка $H/N=K_0\subset K_1\subset\dots\subset K_k=G/N$ с индексами $p$. Прообразы $H_i$ подгрупп $K_i$ в $G$
образуют искомую цепочку: $(H_i:H_{i-1})=(K_i:K_{i-1})=p$.
\end{proof}

\begin{lemma}[3.25]
Пусть $E/F$~--- квадратичное расширение (т.\,е.\ $[E:F]=2$) полей характеристики $\ne2$. Тогда $E=F[\sqrt d]$ для
некоторого $d\in F$.
\end{lemma}

\begin{proof}
Пусть $\alpha\in E\smallsetminus F$. Тогда $F\subsetneq F[\alpha]\subset E$, и по \rf{1.20} $E=F[\alpha]$, так что
минимальный многочлен $\alpha$ имеет степень $2$; пусть это $X^2+bX+c$. Так как $2\ne0$ в $F$, можно выделить полный
квадрат: $(2\alpha+b)^2=b^2-4c$. Значит, $\delta=2\alpha+b$~--- квадратный корень из $d=b^2-4c\in F$, и
$\alpha=\frac{-b+\delta}2$. Поэтому $F[\alpha]=F[\delta]$, т.\,е.\ $E=F[\sqrt d]$.
\end{proof}

В характеристике $2$ лемма неверна. Например, $\F_4=\F_2[\alpha]$, $\alpha^2+\alpha+1=0$, квадратично над $\F_2$, но
$\sqrt0=0$ и $\sqrt1=1$ лежат в $\F_2$, так что $\F_2[\sqrt d]=\F_2$ при любом $d\in\F_2$. Сепарабельные квадратичные
расширения в характеристике $2$~--- это расширения Артина--Шрайера $X^2+X+a$ (\rf{5.32c}); присоединение же
квадратного корня даёт в характеристике $2$ чисто несепарабельное расширение.

\subsection*{Точный критерий построимости}

Условие вещественности поля $E$ в \rf{3.24} не является необходимым, и поэтому \rf{3.24}~--- ещё не критерий.
Чтобы получить критерий, удобно перейти к комплексным числам: отождествим точку $(a,b)$ плоскости с числом $a+bi\in\CC$
и назовём комплексное число $z=a+bi$ \emph{построимым}, если построима точка $(a,b)$. По доказательству \rf{1.36} это
равносильно построимости вещественных чисел $a$ и $b$; вещественное число построимо в старом смысле тогда и только
тогда, когда оно построимо как комплексное число (точка $(a,0)$).

\begin{lemma}[3.25a]
Построимые комплексные числа образуют подполе $\CC$, замкнутое относительно извлечения квадратных корней: если $z$
построимо, то построимы оба числа $w$ с $w^2=z$.
\end{lemma}

\begin{proof}
Пусть $z=a+bi$ и $z'=c+di$ построимы, т.\,е.\ $a,b,c,d$ построимы. Тогда вещественные и мнимые части чисел
\[
z\pm z',\qquad zz'=(ac-bd)+(ad+bc)i,\qquad \frac1z=\frac{a}{a^2+b^2}-\frac{b}{a^2+b^2}\,i\quad(z\ne0)
\]
построимы по \rf{1.37}(a). Пусть теперь $z=a+bi$ построимо, и $r=|z|=\sqrt{a^2+b^2}$; это число построимо по
\rf{1.36}(b) (при $z\ne0$; при $z=0$ доказывать нечего). Так как $r\ge|a|$, числа $\frac{r+a}2$ и $\frac{r-a}2$
неотрицательны, и
\[
x=\sqrt{\tfrac{r+a}2},\qquad y=\pm\sqrt{\tfrac{r-a}2}
\]
построимы. При этом $x^2-y^2=a$ и $(2xy)^2=r^2-a^2=b^2$; выбрав знак $y$ так, чтобы $2xy=b$, получаем $(x+yi)^2=a+bi$.
Второй корень~--- $-(x+yi)$.
\end{proof}

\begin{theorem}[3.25b: критерий построимости]
Для комплексного числа $z$ следующие условия равносильны:
\begin{enumerate}[label=\textup{(\alph*)}]
\item $z$ построимо;
\item $z$ лежит в поле $L_m$ некоторой цепочки подполей $\Q=L_0\subset L_1\subset\dots\subset L_m\subset\CC$ с
$[L_i:L_{i-1}]=2$;
\item $z$ лежит в некотором подполе $E\subset\CC$, которое является расширением Галуа поля $\Q$ степени $2^r$;
\item $z$ алгебраично над $\Q$, и поле разложения его минимального многочлена над $\Q$ имеет степень $2^r$ для
некоторого $r$.
\end{enumerate}
В частности, \emph{вещественное число $\alpha$ построимо тогда и только тогда, когда оно алгебраично и поле разложения
его минимального многочлена имеет над $\Q$ степень, равную степени двойки.} Если $z$ построимо, то между $\Q$ и
$\Q[z]$ есть цепочка полей $\Q=M_0\subset M_1\subset\dots\subset M_s=\Q[z]$ с $[M_j:M_{j-1}]=2$.
\end{theorem}

\begin{proof}
(a)$\Rightarrow$(b). Пусть $z=a+bi$, и $a,b$ построимы. По \rf{1.37}(b) $a\in\Q[\sqrt{a_1},\dots,\sqrt{a_r}]$ и
$b\in\Q[\sqrt{b_1},\dots,\sqrt{b_t}]$, где $a_i$ и $b_j$ положительны и лежат в полях, порождённых предыдущими
корнями. Тогда $a,b$ лежат в поле $K=\Q[\sqrt{a_1},\dots,\sqrt{a_r},\sqrt{b_1},\dots,\sqrt{b_t}]$, а $z\in K[i]$, где
$i=\sqrt{-1}$. На каждом шаге башни $\Q\subset\dots\subset K\subset K[i]$ присоединяется квадратный корень из элемента
предыдущего поля, так что степень шага равна $1$ или $2$; выбросив шаги степени $1$, получаем (b).

(b)$\Rightarrow$(c). По \rf{3.25} $L_i=L_{i-1}[\beta_i]$, $\beta_i^2\in L_{i-1}$, так что $L_m=\Q[\beta_1,\dots,\beta_m]$. Пусть
$\sigma_1=\id,\sigma_2,\dots,\sigma_s$~--- все $\Q$-гомоморфизмы $L_m\to\CC$ (их конечное число, \rf{2.14}), и
\[
E=\Q\bigl[\sigma_j(\beta_i)\,:\,1\le i\le m,\ 1\le j\le s\bigr].
\]
Присоединим образующие по одной в порядке $\sigma_1(\beta_1),\dots,\sigma_1(\beta_m),\sigma_2(\beta_1),\dots,\sigma_2(\beta_m),\dots$.
Квадрат очередного элемента $\sigma_j(\beta_i)$ равен $\sigma_j(\beta_i^2)\in\sigma_j(L_{i-1})=\Q[\sigma_j(\beta_1),\dots,\sigma_j(\beta_{i-1})]$,
т.\,е.\ лежит в уже построенном поле. Значит, каждый шаг имеет степень $1$ или $2$, и $[E:\Q]$~--- степень двойки.
Покажем, что $E/\Q$~--- расширение Галуа. Пусть $\mu_i$~--- минимальный многочлен $\beta_i$ над $\Q$ и
$g=\mu_1\cdots\mu_m$. Каждое $\sigma_j(\beta_i)$~--- корень $\mu_i$. Обратно, если $\gamma$~--- корень $\mu_i$ в $\CC$, то
$\Q$-гомоморфизм $\Q[\beta_i]\to\CC$, $\beta_i\mapsto\gamma$ (\rf{2.1}), продолжается до гомоморфизма $L_m\to\CC$ по \rf{2.12a}
($L_m$ порождено над $\Q[\beta_i]$ корнями многочлена $g$, а $g$ раскладывается в $\CC$ на линейные множители), т.\,е.\
$\gamma=\sigma_j(\beta_i)$ для некоторого $j$. Итак, $E$ порождено над $\Q$ всеми корнями $g$ в $\CC$, т.\,е.\ является полем
разложения $g$; в характеристике $0$ это расширение Галуа (\rf{3.10}). Наконец, $z\in L_m\subset E$.

(c)$\Rightarrow$(d). Элемент $z\in E$ алгебраичен, и так как расширение Галуа $E/\Q$ нормально (\rf{3.10}), $E$ содержит все
корни минимального многочлена $\mu_z$, а значит, и его поле разложения $E_z$. По \rf{1.20} $[E_z:\Q]$ делит
$[E:\Q]=2^r$.

(d)$\Rightarrow$(a). Поле разложения $E_z\subset\CC$ многочлена $\mu_z$~--- расширение Галуа поля $\Q$ (\rf{3.10}), и
$G=\Gal(E_z/\Q)$ имеет порядок $2^r$. Как в доказательстве \rf{3.24} (лемма \rf{3.24a} и основная теорема),
получаем цепочку $E_z=E_0\supset E_1\supset\dots\supset E_r=\Q$ с $E_{i-1}=E_i[\sqrt{d_i}]$, $d_i\in E_i$ (\rf{3.25}). Индукцией по
$i=r,r-1,\dots,0$ покажем, что все элементы $E_i$ построимы. Для $E_r=\Q$ это так (\rf{3.25a}). Если все элементы
$E_i$ построимы, то построимо $d_i$, по \rf{3.25a} построимо $\sqrt{d_i}$, а элементы $E_{i-1}=E_i[\sqrt{d_i}]$~---
многочлены от $\sqrt{d_i}$ с построимыми коэффициентами~--- построимы по \rf{3.25a}. Значит, $z\in E_z$ построимо.

Утверждение о вещественных числах~--- частный случай. Наконец, если $z$ построимо, то в обозначениях (d) подгруппа
$H=\Gal(E_z/\Q[z])$ $2$-группы $G$ по лемме \rf{3.24a} входит в цепочку $H=H_0\subset\dots\subset H_s=G$ с индексами $2$, и
поля $M_j=E_z^{H_{s-j}}$ образуют искомую цепочку между $\Q$ и $E_z^H=\Q[z]$.
\end{proof}

\begin{remark}[3.25c]
(a) Теорема \rf{3.24}~--- частный случай импликации (c)$\Rightarrow$(a) при $E\subset\R$. Вещественности $E$ требовать
нельзя, если мы хотим критерий. Число $\alpha=\sqrt{1+\sqrt2}$ построимо (\rf{1.37}(b) с $a_1=2$, $a_2=1+\sqrt2>0$), но
никакое вещественное расширение Галуа поля $\Q$ его не содержит. Действительно, $[\Q[\alpha]:\Q]=4$: иначе
$\Q[\alpha]=\Q[\sqrt2]$ и $1+\sqrt2=\beta^2$ с $\beta\in\Q[\sqrt2]$; применив автоморфизм $\sqrt2\mapsto-\sqrt2$ поля
$\Q[\sqrt2]$, получили бы $1-\sqrt2=\beta'^2$ с вещественным $\beta'$, что невозможно, так как $1-\sqrt2<0$. Поэтому
минимальный многочлен $\alpha$~--- это $X^4-2X^2-1$ (ведь $\alpha^4=(1+\sqrt2)^2=2\alpha^2+1$). Его корни~--- $\pm\alpha$ и
$\pm\frac i\alpha=\pm i\sqrt{\sqrt2-1}$ (так как $\alpha^2(\sqrt2-1)=1$), и два последних не вещественны. Расширение Галуа,
содержащее $\alpha$, содержит все его сопряжённые (\rf{3.10}), а значит, не лежит в $\R$. Поле разложения~---
$\Q[\alpha,i]$ степени $8$ (так как $i\notin\Q[\alpha]\subset\R$), его группа Галуа~--- диэдральная группа порядка $8$: для
любого корня $\gamma$ автоморфизм $\sigma$ переводит $-\gamma$ в $-\sigma\gamma$, поэтому сохраняет разбиение корней на пары
$\{\pm\alpha\}$, $\{\pm\frac i\alpha\}$, т.\,е.\ лежит в группе симметрий квадрата с вершинами $\alpha,\frac i\alpha,-\alpha,-\frac i\alpha$
(эти пары~--- его диагонали), а она имеет порядок $8=[\Q[\alpha,i]:\Q]$; по
\rf{3.25b} это ещё раз показывает, что $\alpha$ построимо.

(b) Итак, для вещественного $\alpha$ имеем три утверждения:
\begin{itemize}
\item $\alpha$ построимо $\Rightarrow$ $[\Q[\alpha]:\Q]=2^s$ (\rf{1.38});
\item $[\Q[\alpha]:\Q]=2^s$ $\not\Rightarrow$ $\alpha$ построимо (\rf{3.27});
\item $[\Q[\alpha]:\Q]=2^s$ и $\Q[\alpha]/\Q$~--- расширение Галуа $\Rightarrow$ $\alpha$ построимо (\rf{3.24} с $E=\Q[\alpha]\subset\R$).
\end{itemize}
Точный критерий \rf{3.25b} говорит о степени не поля $\Q[\alpha]$, а его нормального замыкания~--- поля разложения
минимального многочлена~$\alpha$.
\end{remark}

\subsection*{Правильные многоугольники}

\begin{corollary}[3.26]
Если $p$~--- простое число вида $2^k+1$, то $\cos\frac{2\pi}p$ построимо.
\end{corollary}

\begin{proof}
Пусть $\zeta=e^{2\pi i/p}$. Поле $\Q[\zeta]$~--- поле разложения $X^p-1$, т.\,е.\ расширение Галуа поля $\Q$, с группой
$G\cong(\Z/p\Z)^\times$ (\rf{5.8}, \rf{1.42}) порядка $p-1=2^k$. Число $\cos\frac{2\pi}p=\frac{\zeta+\zeta^{-1}}2$ лежит в
$\Q[\zeta]$. Группа $G$ абелева, поэтому все её подгруппы нормальны, и по основной теореме (\rf{3.17}) каждое
промежуточное поле, в частности $\Q[\cos\frac{2\pi}p]$, является расширением Галуа поля $\Q$; его степень делит $2^k$
(\rf{1.20}). Так как $\Q[\cos\frac{2\pi}p]\subset\R$, применима теорема \rf{3.24}.
\end{proof}

(Можно и не переходить к вещественному подполю: по \rf{3.25b}(c) построимо уже само $\zeta$, т.\,е.\ точка
$(\cos\frac{2\pi}p,\sin\frac{2\pi}p)$. Точная степень $[\Q[\cos\frac{2\pi}p]:\Q]$ равна $\frac{p-1}2=2^{k-1}$, см.\ рассуждение после \rf{1.42} или \rf{5.11d}.)

Вместе с \rf{1.42} это даёт: \emph{правильный $p$-угольник, $p$ нечётное простое, построим тогда и только тогда, когда
$p$~--- простое число Ферма, т.\,е.\ $p=2^{2^r}+1$} (необходимость доказана в главе 1, достаточность~--- \rf{3.26}; общий
случай правильного $n$-угольника~--- теорема Гаусса \rf{5.12}). Например, мы доказали, что правильный
$65537$-угольник построим, не предъявляя (к счастью) явной формулы для $\cos\frac{2\pi}{65537}$. Доказательство
\rf{3.24}, однако, показывает, \emph{как} такую формулу получить: нужно пройти по цепочке подгрупп группы Галуа. Проделаем
это для $p=17$.

\begin{example}[3.26a: семнадцатиугольник]
Пусть $\zeta=e^{2\pi i/17}$, $G=\Gal(\Q[\zeta]/\Q)\cong(\Z/17\Z)^\times$, $\sigma_a(\zeta)=\zeta^a$. Группа $G$ циклическая
порядка $16$ с образующей $3$: степени $3^0,3^1,\dots,3^{15}$ по модулю $17$ равны
\[
1,3,9,10,13,5,15,11,16,14,8,7,4,12,2,6 .
\]
Поэтому в $G$ ровно по одной подгруппе порядков $8,4,2$ (\rf{5.7b}): $H_8=\langle 9\rangle=\{\pm1,\pm2,\pm4,\pm8\}$
(квадраты), $H_4=\langle13\rangle=\{\pm1,\pm4\}$, $H_2=\{\pm1\}$, и цепочка
$G\supset H_8\supset H_4\supset H_2\supset\{1\}$ отвечает цепочке полей
\[
\Q\subset K_2\subset K_4\subset K_8\subset\Q[\zeta],\qquad [K_d:\Q]=d,\quad K_8=\Q[\zeta]^{\{\pm1\}}=\Q\bigl[\cos\tfrac{2\pi}{17}\bigr]
\]
(\rf{5.11d}). По \rf{5.11e} поле $K_d$ порождено гауссовым периодом~--- суммой $\zeta^a$ по смежному классу
соответствующей подгруппы. Обозначим
\[
\begin{aligned}
&\eta_0=\textstyle\sum_{a\in H_8}\zeta^a,\quad \eta_1=\sum_{a\notin H_8}\zeta^a;\\
&\xi_0=\zeta+\zeta^4+\zeta^{13}+\zeta^{16},\quad \xi_0'=\zeta^2+\zeta^8+\zeta^9+\zeta^{15},\\
&\xi_1=\zeta^3+\zeta^5+\zeta^{12}+\zeta^{14},\quad \xi_1'=\zeta^6+\zeta^7+\zeta^{10}+\zeta^{11};\\
&\lambda=\zeta+\zeta^{16}=2\cos\tfrac{2\pi}{17},\quad \lambda'=\zeta^4+\zeta^{13}=2\cos\tfrac{8\pi}{17}
\end{aligned}
\]
(смежные классы $H_4$~--- это $\{1,4,13,16\}$, $\{2,8,9,15\}$, $\{3,5,12,14\}$, $\{6,7,10,11\}$; первые два составляют $H_8$).
Перемножая суммы и пользуясь равенством $\zeta+\zeta^2+\dots+\zeta^{16}=-1$, получаем
\[
\begin{aligned}
&\eta_0+\eta_1=-1, && \eta_0\eta_1=-4;\\
&\xi_0+\xi_0'=\eta_0, && \xi_0\xi_0'=-1, \qquad \xi_1+\xi_1'=\eta_1,\qquad \xi_1\xi_1'=-1;\\
&\lambda+\lambda'=\xi_0, && \lambda\lambda'=\zeta^5+\zeta^{14}+\zeta^{3}+\zeta^{12}=\xi_1 .
\end{aligned}
\]
(Например, в произведении $\xi_0\xi_0'$ шестнадцать показателей $a+b$, $a\in\{1,4,13,16\}$, $b\in\{2,8,9,15\}$, пробегают
все ненулевые вычеты по модулю $17$ ровно по разу, откуда $\xi_0\xi_0'=-1$; в $\eta_0\eta_1$ каждый ненулевой вычет
встречается $4$ раза.) Итак, на каждом шаге решается квадратное уравнение с коэффициентами из предыдущего поля:
\[
\eta^2+\eta-4=0,\qquad \xi^2-\eta_0\xi-1=0,\qquad \xi^2-\eta_1\xi-1=0,\qquad \lambda^2-\xi_0\lambda+\xi_1=0 .
\]
Знаки корней определяются численными значениями $\eta_0\approx1{,}5616$, $\eta_1\approx-2{,}5616$, $\xi_0\approx2{,}0495$,
$\xi_0'\approx-0{,}4879$, $\xi_1\approx0{,}3442$, $\xi_1'\approx-2{,}9057$, $\lambda\approx1{,}8649$, $\lambda'\approx0{,}1845$:
\[
\eta_{0,1}=\frac{-1\pm\sqrt{17}}2,\qquad \xi_0=\frac{\eta_0+\sqrt{\eta_0^2+4}}2,\qquad
\xi_1=\frac{\eta_1+\sqrt{\eta_1^2+4}}2,\qquad \lambda=\frac{\xi_0+\sqrt{\xi_0^2-4\xi_1}}2 .
\]
Здесь $\lambda=2\cos\frac{2\pi}{17}$. Последний шаг использует $\xi_1$, которое тоже лежит в $K_4=\Q[\xi_0]$ (оно неподвижно относительно $H_4$); в базисе
$1,\xi_0,\xi_0^2,\xi_0^3$ имеем $\xi_1=\frac12(-3+6\xi_0-\xi_0^3)$. Раскрывая эти формулы, получаем выражение
$\cos\frac{2\pi}{17}$ в квадратных радикалах~--- другую запись того же числа, что и в формуле Гаусса из главы~1.
\end{example}

\begin{center}
\begin{keepfig}
\centering
\begin{tikzpicture}[every node/.style={font=\small}]
\node (G) at (0,0) {$G\cong\Z/16$}; \node (H8) at (0,1.1) {$H_8$}; \node (H4) at (0,2.2) {$H_4=\{\pm1,\pm4\}$};
\node (H2) at (0,3.3) {$H_2=\{\pm1\}$}; \node (H1) at (0,4.4) {$\{1\}$};
\foreach \y in {0.55,1.65,2.75,3.85}{\node at (0,\y) {\rotatebox{-90}{$\subset$}};}
\node (Q) at (5.6,0) {$\Q$}; \node (K2) at (5.6,1.1) {$K_2=\Q[\eta_0]=\Q[\sqrt{17}]$};
\node (K4) at (5.6,2.2) {$K_4=\Q[\xi_0]$}; \node (K8) at (5.6,3.3) {$K_8=\Q[\lambda]=\Q[\cos\frac{2\pi}{17}]$};
\node (Z) at (5.6,4.4) {$\Q[\zeta]$};
\foreach \y in {0.55,1.65,2.75,3.85}{\node at (5.6,\y) {\rotatebox{90}{$\subset$}};}
\node[cnew,right] at (8.3,0.55) {$\eta^2+\eta-4=0$};
\node[cnew,right] at (8.3,1.65) {$\xi^2-\eta_0\xi-1=0$};
\node[cnew,right] at (8.3,2.75) {$\lambda^2-\xi_0\lambda+\xi_1=0$};
\node[cnew,right] at (8.3,3.85) {$\zeta^2-\lambda\zeta+1=0$};
\draw[<->,gray] (0.9,0) -- (5.2,0); \draw[<->,gray] (0.5,1.1) -- (3.5,1.1); \draw[<->,gray] (1.3,2.2) -- (4.4,2.2);
\draw[<->,gray] (1.0,3.3) -- (3.3,3.3); \draw[<->,gray] (0.5,4.4) -- (4.9,4.4);
\end{tikzpicture}
\figcap{Пример \rf{3.26a}: цепочка подгрупп циклической группы $(\Z/17\Z)^\times$ и отвечающая ей цепочка полей; каждый шаг
задаётся квадратным уравнением над предыдущим полем.}
\end{keepfig}
\end{center}

\subsection*{Числа степени $4$, которые нельзя построить}

\begin{remark}[3.27]
Обращение к \eqref{eq:c3-constr7} неверно: существуют непостроимые алгебраические числа степени $4$ над $\Q$.

Например, многочлен $f(X)=X^4-4X+2\in\Q[X]$ неприводим (критерий Эйзенштейна при $p=2$), и ниже (\rf{4.10}) мы
покажем, что группа Галуа его поля разложения $E$ равна $S_4$ (его дискриминант $-4864=-2^8\cdot19$ не квадрат, а
кубическая резольвента $X^3-8X-16$ не имеет рациональных корней). У $f$ два вещественных корня ($\approx0{,}5180$ и
$\approx1{,}3631$) и два невещественных; каждый вещественный корень имеет степень $4$. Ни один корень $f$ не построим:
по критерию \rf{3.25b}(d) для этого нужно, чтобы $[E:\Q]$ было степенью двойки, а $[E:\Q]=|S_4|=24$. Полезно
увидеть и прямые рассуждения.

\emph{Первое.} Если бы один корень $f$ был построим, то по \rf{3.25b} построимы были бы все корни (условие (d) зависит
лишь от минимального многочлена), а тогда и все элементы $E$ (\rf{3.25a}: построимые числа образуют поле, а $E$
порождено корнями). Но в $S_4$ есть подгруппа $H$ порядка $8$~--- группа симметрий квадрата с вершинами $1,2,3,4$ (см.\
рисунок; она переставляет вершины квадрата и является силовской $2$-подгруппой $S_4$). Поле $E^H$ имеет степень
$(S_4:H)=3$ над $\Q$, и любой его элемент $z\notin\Q$ имеет степень $3$ (\rf{1.20}). Но построимое $z$ лежит в поле $L_m$ цепочки из
\rf{3.25b}(b), и $[\Q[z]:\Q]$ делит $[L_m:\Q]=2^m$~--- противоречие.

\emph{Второе} (для вещественного корня $\alpha$). Если $\alpha$ построимо, то по \rf{3.25b} есть цепочка
$\Q\subset M\subset\Q[\alpha]$ с $[M:\Q]=[\Q[\alpha]:M]=2$. По основной теореме подгруппы $\Gal(E/M)\supset\Gal(E/\Q[\alpha])$
имеют порядки $12$ и $6$. Подгруппа $\Gal(E/\Q[\alpha])$~--- стабилизатор корня $\alpha$ в $S_4$, т.\,е.\ группа $S_3$
перестановок трёх остальных корней; она содержит транспозиции. А подгруппа $\Gal(E/M)$ индекса $2$ в $S_4$ равна
$A_4$: она нормальна, и факторгруппа порядка $2$ показывает, что вместе с каждым $\sigma$ она содержит $\sigma^2$;
каждый $3$-цикл $\sigma$ есть квадрат ($\sigma=(\sigma^2)^2$), а $3$-циклы порождают $A_4$. В $A_4$ транспозиций нет~---
противоречие. (Иначе говоря, в $A_4$ нет подгруппы порядка $6$: она имела бы индекс $2$ и содержала бы все восемь
$3$-циклов.)

Точно так же не построимы вещественные корни многочлена $X^4-X-1$ из \rf{1.38a} (его группа Галуа тоже $S_4$:
дискриминант $-283$, кубическая резольвента $X^3+4X-1$ неприводима).
\end{remark}

\begin{center}
\begin{keepfig}
\centering
\begin{tikzpicture}[every node/.style={font=\small}]
\begin{scope}[scale=1.1]
\draw[cold,thick] (0,0) -- (1.4,0) -- (1.4,1.4) -- (0,1.4) -- cycle;
\node[below left] at (0,0) {$1$}; \node[below right] at (1.4,0) {$4$}; \node[above right] at (1.4,1.4) {$3$};
\node[above left] at (0,1.4) {$2$};
\draw[cnew,dashed] (0,0) -- (1.4,1.4); \draw[cnew,dashed] (0.7,-0.25) -- (0.7,1.65);
\node[cnew] at (0.7,-0.75) {$H\cong D_4$, $|H|=8$};
\end{scope}
\begin{scope}[xshift=5.2cm,yshift=0.1cm]
\node (Q) at (0,0) {$\Q$}; \node (EH) at (2.2,0) {$E^H$}; \node (E) at (4.4,0) {$E$};
\node at (1.1,0) {$\subset$}; \node at (3.3,0) {$\subset$};
\draw[decorate,decoration={brace,amplitude=5pt}] (0,0.3) -- (2.2,0.3) node[midway,above=5pt,cnew] {$3$};
\draw[decorate,decoration={brace,amplitude=5pt}] (2.2,0.3) -- (4.4,0.3) node[midway,above=5pt] {$8$};
\draw[decorate,decoration={brace,mirror,amplitude=5pt}] (0,-0.3) -- (4.4,-0.3) node[midway,below=5pt,cold] {$24$};
\end{scope}
\end{tikzpicture}
\figcap{Слева: подгруппа $H\subset S_4$ порядка $8$~--- симметрии квадрата (повороты и отражения, переставляющие
вершины $1,2,3,4$). Справа: её неподвижное поле $E^H$ имеет степень $3$ над $\Q$, поэтому его элементы вне $\Q$ не
построимы.}
\end{keepfig}
\end{center}
